\chapter{Hyperkalemia (High Potassium)}

\section*{Definition}
Hyperkalemia is a higher-than-normal concentration of potassium in the blood. It can impair cardiac conduction and cause life-threatening arrhythmias, particularly when the increase is rapid or substantial.

\section*{When to See a Doctor}
Severe weakness, paralysis, palpitations, chest pain, fainting, or a reported high potassium level requires urgent medical assessment. Do not attempt to correct a significant result with home remedies.

\section*{How It Develops}
Potassium establishes the electrical gradient across cell membranes. Excess extracellular potassium partially depolarizes cardiac and skeletal muscle cells, producing weakness and characteristic ECG changes that can progress to conduction block, ventricular arrhythmia, or cardiac arrest.

\section*{Causes}
\begin{itemize}
  \item Acute or chronic kidney disease, adrenal insufficiency, dehydration, and severe tissue breakdown.
  \item Medicines such as ACE inhibitors, ARBs, potassium-sparing diuretics, NSAIDs, trimethoprim, and potassium supplements.
  \item Metabolic acidosis, uncontrolled diabetes, rhabdomyolysis, hemolysis, tumor lysis, and transfusion-related potassium load.
  \item Pseudohyperkalemia from difficult blood collection, prolonged tourniquet use, or sample hemolysis.
\end{itemize}

\section*{Related Symptoms}
\begin{itemize}
  \item Often asymptomatic; weakness, paresthesias, nausea, palpitations, or shortness of breath may occur.
  \item ECG abnormalities may include peaked T waves, PR prolongation, QRS widening, or a sine-wave pattern.
\end{itemize}

\section*{Medical Evaluation}
\begin{itemize}
  \item Repeat potassium promptly if pseudohyperkalemia is possible, with renal function, bicarbonate, glucose, and medication review.
  \item Obtain an ECG urgently when hyperkalemia is moderate, severe, symptomatic, rapidly rising, or associated with kidney disease.
  \item Investigate kidney function, acid-base status, endocrine disease, tissue injury, and medication or supplement exposure.
\end{itemize}

\section*{Treatment Options}
Emergency treatment may include intravenous calcium to stabilize the myocardium, insulin with glucose, inhaled beta-agonist, bicarbonate in selected acidotic patients, potassium-binding therapy, diuresis, and dialysis. The cause must be corrected and potassium rechecked.

\section*{Self-Care and Home Management}
Contact the treating clinician promptly about an abnormal result. Do not take potassium supplements or potassium-containing salt substitutes unless prescribed. Do not discontinue heart or blood-pressure medicine without medical advice.

\section*{References}
\begin{thebibliography}{99}
\bibitem{hyperkalemia:ref1} Mount DB. Disorders of potassium balance. In: \textit{Harrison's Principles of Internal Medicine}. 21st ed. McGraw-Hill; 2022.
\bibitem{hyperkalemia:ref2} UK Kidney Association. Clinical practice guideline: treatment of acute hyperkalaemia in adults. UK Kidney Association; 2023.
\end{thebibliography}
